\label{cap:introducao_modelo}

O presente Trabalho de Conclusão de Curso foi desenvolvido em cooperação técnico-científica, em coautoria com o acadêmico Victor de Souza Vilela da Silva, do Curso de Engenharia de Computação do Centro Federal de Educação Tecnológica de Minas Gerais, Unidade Leopoldina. Esta pesquisa faz parte de um estudo complementar e interdependente intitulado \textit{Interpretador Dinâmico para Ensino de Programação: Implementação de análise léxica e sintática, geração e interpretação de código intermediário em ambiente web}, de autoria do referido coautor e submetido ao mesmo curso desta Instituição. A divisão das atividades comuns e dos eixos individuais de aprofundamento é descrita no Apêndice~\ref{ap:plano}.

Embora este texto e o documento complementar acima identificado compartilhem a mesma base teórica geral e o mesmo objeto de estudo, isto é, uma plataforma educacional web formada por um interpretador dinâmico integrado a uma IDE com gramática configurável pelo usuário, as abordagens diferem em escopo e ênfase metodológica. Enquanto o documento complementar, de autoria de Victor de Souza Vilela da Silva, se concentra no núcleo de tradução, examinando as fases de análise léxica, análise sintática, geração de código intermediário e execução pelo interpretador, este texto trata do Ambiente de Desenvolvimento Integrado: sua arquitetura \textit{front-end}, a criação e implementação do assistente de personalização da linguagem, a experiência de uso da plataforma e a preparação de uma avaliação futura com estudantes. Para uma compreensão completa do sistema proposto, recomenda-se a leitura dos dois documentos.

O ensino introdutório de programação envolve dificuldades de compreensão e aplicação de conceitos, discutidas no mapeamento sistemático de \citeonline{souza2016}. Este trabalho enfoca uma dessas dificuldades: lidar simultaneamente com a construção do algoritmo e com as convenções de escrita da linguagem. Os experimentos de \citeonline{stefik2013} indicam que escolhas de sintaxe afetam o desempenho de iniciantes, motivando investigar outras formas de apresentação dos comandos. Esse resultado não demonstra, por si só, que a personalização proposta neste projeto reduza reprovação, evasão ou carga cognitiva.

Neste trabalho, \textbf{dinâmico} designa a configuração recebida pelo núcleo a cada nova análise e execução: palavras-chave, operadores em forma de palavra, literais booleanos e delimitadores podem ser redefinidos, e o usuário pode selecionar modos de terminação de instruções, blocos, tipagem e arranjos já implementados. As escolhas são aplicadas sem recompilar a aplicação web; não permitem acrescentar produções gramaticais arbitrárias nem mudar a semântica das operações durante a execução de um programa. A Tabela~\ref{tab:elementos-configuraveis} delimita essas possibilidades. O núcleo compila o código-fonte para instruções intermediárias e as executa por interpretação; a IDE oferece a interface para configurar e acionar esse processo.

A pergunta que orienta esta pesquisa nasce dessa observação: \textbf{em que medida uma plataforma de programação cuja sintaxe possa ser redefinida pelo próprio estudante contribui para reduzir a barreira sintática enfrentada no primeiro contato com a programação, deslocando o foco da forma para o conteúdo do raciocínio computacional?}

Trabalha-se com a hipótese, ainda não submetida a verificação empírica, de que permitir ao estudante escolher quais palavras servirão como comandos, como os blocos serão delimitados, como os operadores serão escritos e se o terminador de instrução será obrigatório ou opcional possa reduzir a carga cognitiva exigida no aprendizado simultâneo da lógica e da forma. A expectativa é que o exercício de adaptar a linguagem ao próprio repertório converta a sintaxe, antes vivida como obstáculo, em objeto consciente de escolha pedagógica.

Cumpre delimitar, desde já, o alcance deste trabalho diante dessa pergunta. O que aqui se apresenta é a construção do artefato que torna a investigação possível e uma proposta de protocolo com que ela poderá ser conduzida. A aplicação desse protocolo junto a estudantes, que permitiria responder empiricamente à pergunta formulada, não integra o escopo deste documento e é registrada como encaminhamento futuro.

\section{Justificativa}

A justificativa parte de uma comparação delimitada entre ferramentas. O Portugol Studio utiliza comandos em português e oferece depuração e inspeção de variáveis \cite{univali2026}. O Quorum investiga decisões de linguagem com base em evidências e oferece recursos de acessibilidade \cite{quorum2025}. O Beecrowd apoia a prática por submissões avaliadas automaticamente \cite{beecrowd2026faq}. Essas abordagens já tratam dificuldades de iniciantes, mas não têm como foco o mesmo assistente de configuração de vocabulário e de variantes sintáticas proposto aqui.

A contribuição deste trabalho é integrar esse assistente a uma IDE web com editor, terminal, visualização de \textit{tokens} e código intermediário, depuração e recursos de gestão de exercícios. A comparação da Seção~\ref{sec:trabalhos-relacionados} explicita os critérios e os limites desse recorte, sem generalizar seus resultados para todos os ambientes educacionais. A efetividade pedagógica da combinação permanece uma hipótese, cuja avaliação requer dados de uso e aprendizagem.

\section{Objetivos}

A seguir são apresentados o objetivo geral que orienta a investigação e os objetivos específicos que viabilizam o seu alcance.

\subsection{Objetivo Geral}

Projetar, implementar e integrar uma IDE web ao núcleo de tradução e execução configurável, oferecendo ao estudante iniciante recursos para personalizar a linguagem e experimentar algoritmos no navegador.

\subsection{Objetivos Específicos}

Para que o objetivo geral seja efetivamente alcançado, as seguintes metas operacionais foram fixadas:

\begin{itemize}
    \item Especificar quais elementos da linguagem serão passíveis de redefinição pelo usuário, delimitando o conjunto configurável de palavras-chave, operadores, literais, delimitadores de bloco e terminador de instrução;
    \item Projetar e implementar a IDE web, contemplando o assistente de personalização da linguagem, a edição de código, a execução no próprio ambiente e a apresentação de \textit{feedback} ao usuário;
    \item Integrar a IDE ao interpretador dinâmico, de modo que as análises léxica e sintática e a execução do código intermediário sejam acionadas a partir do código produzido no editor;
    \item Disponibilizar a plataforma em ambiente web, com os fluxos de uso concluídos e com atenção a requisitos de acessibilidade e responsividade;
    \item Propor um protocolo de avaliação da plataforma, estabelecendo os instrumentos de coleta, as métricas e os critérios que permitirão, em etapa posterior, aferir sua contribuição pedagógica.
\end{itemize}

\section{Estrutura do Trabalho}

O restante deste trabalho está organizado da seguinte maneira. O Capítulo~\ref{cap:referencial} apresenta o referencial teórico que sustenta o projeto, reunindo os fundamentos de construção de compiladores, as práticas associadas ao desenvolvimento de aplicações web e as iniciativas educacionais relacionadas. O Capítulo~\ref{cap:proposta}, por sua vez, expõe a proposta do trabalho e o planejamento de suas duas fases. Por fim, o Capítulo~\ref{cap:metodologia} descreve a modelagem realizada e a implementação consolidada até o encerramento desta etapa, contemplando a arquitetura da plataforma, a personalização dinâmica da linguagem, a integração com o interpretador e a estratégia de testes adotada.

O Capítulo~\ref{cap:consideracoes} retoma os objetivos, apresenta os resultados técnicos e as limitações e delimita os trabalhos futuros. A Seção~\ref{sec:protocolo-avaliacao} apresenta a proposta de avaliação, ainda não aplicada com estudantes.
